% Responsable: P2
\section{Implementación del parser}
\label{sec:parser}

\subsection{Un método por no terminal}

El parser (\texttt{spanish/parser.py}) es descendente recursivo: cada no terminal de la
gramática sin recursión por la izquierda es un método de la clase \texttt{Parser}, que consume
las palabras que le corresponden y devuelve el nodo del árbol, o \texttt{None} si no puede.

\begin{table}[h]
  \centering
  \small
  \resizebox{\linewidth}{!}{%
  \begin{tabular}{lll}
    \toprule
    No terminal & Método & Cómo se programa \\
    \midrule
    $S \to NP\ VP$ & \texttt{sentence} & secuencia \\
    $NP \to NPbase\ NP'$ & \texttt{noun\_phrase} & \texttt{while} para cada $PP$ ($NP'$) \\
    $NPbase \to \textit{Prn} \mid \textit{Det}\ Nom \mid \textit{Pos}\ Nom$ & \texttt{noun\_phrase\_base} & elección por la palabra actual \\
    $Nom \to \textit{Noun}\ AdjOpc$ & \texttt{nominal} & \texttt{if} para el $AdjP$ opcional \\
    $AdjP \to \textit{Adj} \mid \textit{Intens}\ \textit{Adj}$ & \texttt{adjective\_phrase} & intensificador opcional \\
    $VP \to \textit{Verb}\ ObjOpc\ VP'$ & \texttt{verb\_phrase} & \texttt{if} para el objeto, \texttt{while} para $PP$ y \textit{Adv} \\
    $PP \to \textit{Prep}\ NP$ & \texttt{prepositional\_phrase} & secuencia \\
    \bottomrule
  \end{tabular}}
  \caption{Correspondencia entre la gramática y el código.}
  \label{tab:metodos_parser}
\end{table}

Tres operaciones sostienen todos los métodos:
\begin{itemize}
  \item \texttt{starts(X)} dice si la palabra actual puede empezar el no terminal $X$, según su
    conjunto FIRST (\texttt{spanish/grammar.py}). Así se programan las partes opcionales
    ($\varepsilon$) y las repeticiones.
  \item \texttt{match(c)} consume la palabra si tiene una lectura de la categoría $c$.
  \item \texttt{expect(c, descripción)} es igual, pero si la categoría falta registra el error
    con la descripción de lo que se esperaba.
\end{itemize}
Una oración se acepta solo si \texttt{sentence} construye el árbol \emph{y} no sobran palabras.
El árbol se arma ya en el formato acordado: $NP'$, $AdjOpc$, $ObjOpc$ y $VP'$ no aparecen, y sus
hijos cuelgan del nodo padre (sección~\ref{sec:gramatica}).

\subsection{Palabras ambiguas}

El lexer no decide la categoría de una palabra: le entrega al parser todas las lecturas que
tiene en el léxico. Hay cinco palabras con más de una: \emph{bajo} (adjetivo o preposición) y
cuatro formas verbales que son iguales en presente y en pretérito (\emph{matamos},
\emph{dormimos}, \emph{compramos}, \emph{buscamos}). El parser elige la lectura que pide la
posición en la que está: en \emph{El gato duerme bajo la mesa}, después del verbo solo puede
empezar un $PP$, así que \emph{bajo} se lee como preposición; en \emph{El hombre bajo come},
después del sustantivo se lee como adjetivo.

Como la gramática es LL(1) salvo por la frase preposicional, esta estrategia funciona sin
retroceso en todo el conjunto de prueba. Pero tiene un límite: en \emph{El gato bajo la mesa
duerme}, que la gramática sí genera (\emph{bajo la mesa} como frase preposicional del sujeto), el
parser toma \emph{bajo} como adjetivo porque puede empezar un $AdjP$, y luego falla en
\emph{la}. La palabra ambigua vuelve a poner \textit{Prep} en FIRST($AdjOpc$) y en
FOLLOW($AdjOpc$). Resolverlo requiere retroceso (probar la otra lectura al fallar) o un parser
que explore todas las alternativas, como un parser de gráfico (\emph{chart parser}).

Las formas verbales ambiguas no causan ese problema: sus dos lecturas son de la misma categoría
y solo cambia el tiempo, que no interviene en la concordancia.

\subsection{Pipeline y mensajes de error}

\texttt{spanish/pipeline.py} encadena las cuatro etapas, igual que en la Fase~1: si una etapa
falla, la siguiente no se ejecuta.
\begin{center}
  texto $\to$ \textbf{lexer} $\to$ palabras con lecturas $\to$ \textbf{parser} $\to$ árbol
  $\to$ \textbf{concordancia} $\to$ aceptada o rechazada $\to$ \textbf{visualización}
\end{center}
Cada error dice en qué palabra ocurrió y por qué. El parser reporta solo el primer error
sintáctico, porque después de él el resto del análisis ya no tiene sentido; el analizador de
concordancia los reporta todos (tabla~\ref{tab:errores_demo}).

\begin{table}[h]
  \centering
  \small
  \begin{tabular}{p{0.3\linewidth}p{0.62\linewidth}}
    \toprule
    Oración & Salida del pipeline \\
    \midrule
    \emph{Ella come una manzana\#} & error léxico: carácter inválido `\#' (columna 22) \\
    \emph{El perro corre en el parque.} & error sintáctico, palabra 3: `corre' no está en la gramática \\
    \emph{Come la manzana.} & error sintáctico, palabra 1: se esperaba un pronombre, determinante o posesivo \\
    \emph{El gato come muy.} & error sintáctico, palabra 4: no se esperaba `muy' \\
    \emph{Los perros come.} & error de concordancia, palabra 3: el verbo `come' (singular) no concuerda con el sujeto `perros' (plural) \\
    \emph{La elefante grises duerme.} & dos errores de concordancia: el determinante `La' no concuerda en género con `elefante'; el adjetivo `grises' no concuerda en número con `elefante' \\
    \bottomrule
  \end{tabular}
  \caption{Mensajes de error del pipeline.}
  \label{tab:errores_demo}
\end{table}

\subsection{Uso}

El pipeline se usa desde la terminal, para la demo, o desde Python, para los experimentos:
\begin{lstlisting}
$ python -m spanish.pipeline "Él mató un elefante muy gris ayer." --arbol
$ python -m spanish.pipeline "Los perros come." --tokens
$ python -m evaluation.run_parser        # las 73 oraciones -> results/parser.jsonl
\end{lstlisting}
Además de aceptar o rechazar, el resultado incluye las etiquetas de cada palabra, el núcleo del
sujeto, el verbo, el núcleo del objeto y a qué se une la primera frase preposicional. Esos datos
salen del árbol, así que existen cuando la oración se pudo derivar, aunque la rechace la
concordancia. Son los que se comparan con los demás sistemas en la
sección~\ref{sec:resultados}.

\subsection{Pruebas}

El repositorio tiene 220 pruebas automáticas (\texttt{pytest}). Las del parser y el pipeline
verifican la gramática, el formato de los árboles y los mensajes de error, y pasan las 73
oraciones del conjunto de prueba por el pipeline completo. Esa prueba asegura que lo que el
conjunto dice que la gramática acepta coincide con lo que el parser acepta.
